%! Author = zhiweizhang
%! Date = 10.06.25

\chapter{Conclusion and Future Work}\label{chapter:conclusion-and-future-work}

\section{Conclusion}

This thesis presented a comprehensive study on the classification of Nazi symbols in images using a range of models,
including convolutional neural networks (\ac{ResNet}, EfficientNet), transformer-based architectures (\ac{ViT}), object
detection models (YOLOv11), and multimodal models (OpenCLIP). In addition to evaluating end-to-end deep learning
approaches, we also explored hybrid pipelines that combined OpenCLIP image embeddings with classical machine
learning classifiers.

The models were evaluated on two classification tasks: binary classification (Nazi vs. non-Nazi symbols) and
multi-class classification (discriminating among 13 specific Nazi-related symbol classes and a non-Nazi class). For
the binary task, the best-performing models were OpenCLIP combined with classical classifiers (e.g., \ac{SVC} and \ac{KNN}),
\ac{ViT}, and YOLOv11, all achieving F1-scores above 0.89 and \ac{AUC} values above 0.94. This indicates that both
learned embeddings and high-resolution feature extractors are effective for distinguishing Nazi symbols from benign content.

In the multi-class setting, YOLOv11 emerged as the top performer, followed closely by \ac{ViT}. OpenCLIP alone performed
poorly in zero-shot settings, but its performance improved substantially when paired with supervised
classifiers—though still underperforming compared to end-to-end deep models. EfficientNet showed weak results in
both settings, suggesting sensitivity to architectural choices and training hyperparameters.

Overall, our results demonstrate that symbol classification, especially in sensitive and subtle contexts like
detecting hate symbols, benefits from models that are either end-to-end optimized for the specific task or from
hybrid approaches that combine semantic embeddings with suitable classifiers. The experiments also highlight the
challenges in achieving reliable generalization across diverse classes of hate symbols, given their visual
variability and contextual ambiguity.

\section{Future Work}

This study has demonstrated the feasibility of using deep learning and multimodal models to detect Nazi and
extremist symbols. Building on these results, several directions for future work emerge:

\begin{itemize}
    \item \textbf{Dataset Expansion and Annotation:} Enhancing the dataset—particularly increasing the number of
    examples for underrepresented or visually similar symbol classes—could significantly improve multi-class
    classification performance. Incorporating domain expert feedback or leveraging semi-supervised and weakly
    supervised learning strategies could further refine annotation quality and contextual relevance.

    \item \textbf{Fine-Tuning Multimodal Models:} While OpenCLIP embeddings combined with classical classifiers
    showed promise, future work should explore fine-tuning the vision backbone or jointly training multimodal models
    on task-specific data. This may improve class separability and allow better adaptation to symbol-specific nuances.

    \item \textbf{Explainability and Robustness:} The integration of \ac{XAI} techniques (e.g.,
    \ac{Grad-CAM}, \ac{SHAP}) would help interpret classification decisions—especially important in sensitive applications
    such as legal review or moderation. Evaluating model robustness to adversarial attacks, image distortions, or
    low-quality inputs remains an important next step.

    \item \textbf{Context-Aware and Multimodal Classification:} Incorporating metadata such as image captions,
    filenames, or textual context from surrounding content may improve performance, particularly in ambiguous cases.
    This aligns with the broader goal of context-aware content moderation in social platforms and digital archives.

    \item \textbf{Online Deployment and Real-Time Integration:} Although this work implemented and containerized an
    \ac{API} service for local deployment, future work could involve deploying it in an online setting using cloud
    platforms. This includes optimizing for inference speed, \ac{API} scalability, and data privacy compliance in
    real-time systems.

    \item \textbf{User Feedback and Human-in-the-Loop Systems:} Integrating user feedback into the classification
    pipeline could help refine model outputs over time and improve trustworthiness. This is especially relevant in
    professional or archival environments, where user corrections can guide continuous learning.

\end{itemize}

In summary, this thesis lays a foundation for scalable and practical Nazi symbol classification. Extending it to
real-world systems will require continued efforts in model refinement, dataset curation, deployment optimization,
and ethical safeguards.

